\begin{abstract}
多变量时间序列缺失修复需要从时间与变量维度上分布不均的观测中重建局部轨迹。其关键在于融合互补的上下文与Patch信息，并保持恢复数值与真实可见测量的一致性。本文提出证据引导的粗到精修复框架\Model{}，包含时间--变量上下文编码、多来源Patch证据检索、逐点证据引导重建和观测条件化残差精恢复四个组件。上下文编码捕捉时间动态与跨变量依赖；Patch检索将当前窗口及训练历史片段组织为统一证据库，由各目标点分别读取以增强上下文基础估计；残差精恢复进一步结合初始估计、可见观测和观测导出的参考信息，纠正重建误差。实验覆盖六个不同应用领域的数据集。在与容量匹配上下文模型的三训练种子配对实验中，证据路径使ETTh1和PAMAP2的最终MSE分别降低\GainETTMSE{}和\GainPAMMSE{}。组件、缺失条件及效率分析进一步呈现Patch证据与观测条件化精恢复对缺失重建的作用。
\end{abstract}
